\section{Related Work}
\label{sec:related-work}

\subsection{Garment manipulation}

Garment manipulation has commonly been made tractable by reducing control to
task-specific primitives, such as selecting grasp points followed by a fixed
pick--place or fling motion. These methods encode useful geometric and control
knowledge and can solve well-defined garment tasks efficiently. At the same
time, the primitive determines much of the behavior before learning begins. It
therefore provides limited evidence about whether a generalist policy can infer
garment state, produce continuous robot motion, and recover from its own
mistakes. ClothDojo retains the ordinary robot action interface so that these
abilities remain part of the evaluated problem.

\subsection{Deformable simulation and benchmarks}

Simulation environments, simulation benchmarks, and policy benchmarks answer
different questions. An environment determines what can be simulated; a
fidelity benchmark measures how well simulated cloth matches physical cloth;
a policy benchmark measures what an agent can accomplish under a fixed task
distribution. GarmentLab and RGBench primarily strengthen the first two
layers. SimWeaver additionally connects simulation to policy training and is
therefore the closest setting to ClothDojo. ClothDojo concentrates on the third
layer: severe garment flattening with continuous actions, scalable successful
data, and closed-loop evaluation across garment variation.

\subsection{Data generation for generalist robot policies}

Privileged supervision separates the difficulty of producing useful training
behavior from the difficulty of the final visual task. NOCS supplies stable
cross-asset geometric cues to the trajectory agent while preserving the VLA's
image-to-action structure. Once a successful physical trajectory exists, its
observations can be rendered without those cues. Thus the downstream policy
still acts from ordinary images and never accesses particle state; privilege is
used to manufacture data, not to simplify the evaluated policy.
